% Part: model-theory
% Chapter: basics
% Section: reducts-and-expansions

\documentclass[../../../include/open-logic-section]{subfiles}

\begin{document}

\olfileid{mod}{bas}{red}
\section{ન્યૂનરૂપો અને વિસ્તારો}

સામાન્ય સંકેતો ધરાવતી ભાષાઓની અને એ ભાષાઓ માટેની !!{structure}sની
સરખામણી કરવી ઘણી વાર ઉપયોગી અથવા જરૂરી હોય છે. સૌથી સામાન્ય કિસ્સામાં
!!a{language}~$\Lang{L}$ના બધા સંકેતો !!a{language}~$\Lang{L'}$ના પણ
ભાગ હોય છે, એટલે કે $\Lang{L} \subseteq \Lang{L'}$. ત્યારે વધારાના
સંકેતોનાં અર્થઘટનો ઉમેરીને અને સામાન્ય સંકેતોનાં અર્થઘટનો યથાવત રાખીને
$\Lang{L}$-!!{structure}~$\Struct{M}$ને હંમેશાં
$\Lang{L'}$-!!{structure}માં વિસ્તારી શકાય છે. બીજી તરફ,
$\Lang{L'}$-સંરચના~$\Struct{M'}$માંથી~$\Lang{L}$માં ન આવતા સંકેતોનાં
અર્થઘટનોને માત્ર ``ભૂલી જઈને'' આપણે $\Lang{L}$-સંરચના મેળવી શકીએ છીએ.

\begin{defn}
\ollabel{defn:reduct}
ધારો કે $\Lang L \subseteq \Lang L'$, $\Struct M$ એ
$\Lang L$-!!{structure} છે અને $\Struct M'$ એ
$\Lang L'$-!!{structure} છે. નીચેની શરતો સંતોષાય તો અને માત્ર તો
$\Struct M$ને $\Lang L$ પરનું $\Struct M'$નું \emph{ન્યૂનરૂપ} અને
$\Struct M'$ને $\Lang L'$ સુધીનો $\Struct M$નો \emph{વિસ્તાર} કહીએ:
\begin{enumerate}
\item $\Domain{M} = \Domain{M'}$
\item દરેક !!{constant}~$c \in \Lang L$ માટે $\Assign{c}{M} =
  \Assign{c}{M'}$.
\item દરેક !!{function}~$f \in \Lang L$ માટે $\Assign{f}{M} =
  \Assign{f}{M'}$.
\item દરેક !!{predicate}~$P \in \Lang L$ માટે $\Assign{P}{M} =
  \Assign{P}{M'}$.
\end{enumerate}
\end{defn}

\begin{prop}
\ollabel{prop:reduct}
જો $\Lang{L}$-!!{structure}~$\Struct{M}$ એ
$\Lang{L'}$-!!{structure}~$\Struct{M'}$નું ન્યૂનરૂપ હોય, તો બધા
$\Lang{L}$-!!{sentence}s~$!A$ માટે
\[
\Sat{M}{!A} \text{ ત્યારે અને માત્ર ત્યારે, જ્યારે } \Sat{M'}{!A}.
\]
\end{prop}

\begin{proof}
  અભ્યાસપ્રશ્ન.
\end{proof}

\begin{prob}
\olref[mod][bas][red]{prop:reduct} સાબિત કરો.
\end{prob}

\begin{defn}
આપણી પાસે $\Lang{L}$-સંરચના $\Struct{M}$ હોય, એક જ $n$-સ્થાની
!!{predicate}~$P$ ઉમેરીને મળતો $\Lang{L}$નો વિસ્તાર $\Lang{L'} =
\Lang{L} \cup \{P\}$ હોય અને $R \subseteq \Domain{M}^n$ એ
$n$-સ્થાની સંબંધ હોય, ત્યારે $\Assign{P}{M'} = R$ ધરાવતા
$\Struct{M}$ના વિસ્તાર~$\Struct{M'}$ માટે આપણે $\Expan{M}{R}$ લખીએ છીએ.
\end{defn}

\end{document}
